\documentclass[12pt, a4paper]{article}
\usepackage[utf8]{inputenc}
\usepackage[T1]{fontenc}
\usepackage{amsmath, amssymb, amsthm}
\usepackage{geometry}
\usepackage{hyperref}
\geometry{margin=2.5cm}

\newtheorem{lema}{Lema}[section]
\newtheorem{corolario}[lema]{Corolário}
\newtheorem{proposicao}[lema]{Proposição}
\newtheorem{observacao}[lema]{Observação}
\newtheorem{definicao}[lema]{Definição}
\newtheorem{exemplo}[lema]{Exemplo}

\title{Uma propriedade de soma constante na grade $3\times N$:\\
	Invariante Âncora, estrutura aritmética e conexões\\
	com as conjecturas de Goldbach}
\author{Tiago Bandeira\\
	\small\texttt{tiagobandeira.goldbach2025@gmail.com}}
\date{Maio de 2026\\[4pt]
	\small Preprint — Paper 2 da série; complementa \cite{bandeira2026}}

\begin{document}
	
	\maketitle
	
	\begin{abstract}
		A grade $G_{3,N}$, formada pelos primeiros $3N$ números ímpares organizados em três
		linhas, possui uma propriedade aritmética notável: em cada janela $3\times 3$ as três
		configurações canônicas (diagonal principal $D_1$, diagonal secundária $D_2$ e coluna
		central $C_v$) têm sempre a mesma soma $S(i)=3(2N+2i+1)$. Demonstramos esse
		\textbf{Lema da Soma Constante} de forma rigorosa e derivamos, como consequência
		imediata, o \textbf{Invariante Âncora}: o elemento central da linha mediana satisfaz
		$f(2,i{+}1)=S(i)/3$, determinando por sua posição o \textit{complemento aritmético}
		oracle-free $C(n,i)=S(i)-n$ de qualquer outro elemento da janela.
		
		Mostramos que cada configuração canônica é uma progressão aritmética com razão
		determinada por $N$, e que as diferenças entre linhas são constantes — revelando a
		estrutura modular uniforme dos elementos da grade. Generalizamos o resultado para
		grades $k\times N$: a soma da coluna central é sempre $k$ vezes o elemento da linha
		mediana quando $k$ é ímpar; para $k$ par não existe âncora inteira central, e para
		$k=4$ o complemento $C$ é sempre par mas a âncora é racional.
		
		Finalmente, mostramos que quando uma configuração é totalmente prima ela produz
		simultaneamente três representações da Conjectura de Goldbach binária (forte) e uma
		da Conjectura de Goldbach ternária (fraca), e que o complemento $C(n,i)$ é sempre
		par — fato que, explorado em \cite{bandeira2026c}, fundamenta o estimador
		geométrico oracle-free $\Omega^*_N$.
	\end{abstract}
	
	\tableofcontents
	\bigskip
	
	%----------------------------------------------------------------------
	\section{Introdução}
	%----------------------------------------------------------------------
	
	A grade $G_{3,N}$ é construída organizando os primeiros $3N$ números ímpares
	consecutivos em três linhas de $N$ colunas. Introduzida em~\cite{bandeira2026} para
	estudar a distribuição de primos em configurações geométricas rígidas, ela revelou
	propriedades aritméticas que não eram aparentes na reta numérica usual. Em particular,
	durante a investigação do \textit{peso geométrico} $\Omega_N(n)$ — que mede a fração
	de vizinhos primos de uma célula — observou-se uma identidade algébrica simples mas
	central: em qualquer submatriz $3\times 3$ (denominada \textbf{janela} $W_i$), as três
	configurações canônicas têm exatamente a mesma soma.
	
	Este artigo apresenta essa propriedade de forma independente e autocontida, e
	desenvolve todas as consequências que ela implica para a série de trabalhos sobre a
	Conjectura de Goldbach.
	
	\medskip
	
	O papel deste artigo no contexto da série é estrutural: ele fornece a fundação
	algébrica que torna possível o passo seguinte. O peso geométrico $\Omega_N(n)$
	introduzido em~\cite{bandeira2026} depende de um \textit{oráculo de primalidade} —
	é preciso saber quais vizinhos são primos para calculá-lo. Isso impede sua utilização
	como crivo \textit{a priori}. O Invariante Âncora (Seção~\ref{sec:ancora}) resolve
	esse obstáculo: como $f(2,i{+}1)=S(i)/3$ é determinado apenas pela posição na grade, o
	complemento $C(n,i)=S(i)-n$ é computável sem qualquer teste de primalidade. Com ele,
	o trabalho~\cite{bandeira2026c} constrói o estimador oracle-free $\Omega^*_N$, que
	supera $\Omega_N$ em robustez computacional.
	
	\medskip
	
	As principais contribuições deste artigo são:
	\begin{itemize}
		\item O \textbf{Lema da Soma Constante} (Seção~\ref{sec:lema}): $D_1=D_2=C_v=S(i)$
		para toda janela $W_i$.
		\item O \textbf{Invariante Âncora} (Seção~\ref{sec:ancora}): $f(2,i{+}1)=S(i)/3$
		e a paridade garantida do complemento $C(n,i)$.
		\item A \textbf{estrutura aritmética das configurações} (Seção~\ref{sec:aritm}):
		cada configuração canônica é uma progressão aritmética com razão que depende
		de $N$; diferenças entre linhas são constantes; estrutura modular uniforme.
		\item A \textbf{generalização para grades $k\times N$} (Seção~\ref{sec:general}):
		análise da âncora para $k$ arbitrário, com distinção entre $k$ par e ímpar.
		\item As \textbf{conexões com Goldbach} (Seção~\ref{sec:goldbach}): cada trio primo
		produz três representações binárias e uma ternária; o complemento $C(n,i)$
		como fundamento oracle-free.
	\end{itemize}
	
	\medskip
	
	Enquanto a Conjectura de Goldbach fraca (todo ímpar $>5$ é soma de três primos) foi
	demonstrada por Helfgott~\cite{helfgott2013}, e o resultado de Vinogradov cobre os
	inteiros suficientemente grandes, a conexão geométrica aqui apresentada tem valor
	estrutural independente: ela exibe, no mesmo objeto algébrico, testemunhas das duas
	formas da conjectura, e fornece o complemento $C$ como novo objeto de estudo cujas
	propriedades aditivas são centrais ao programa do crivo geométrico.
	
	%----------------------------------------------------------------------
	\section{Definição da grade $G_{3,N}$ e das janelas}
	\label{sec:definicao}
	%----------------------------------------------------------------------
	
	\subsection{Grade $G_{3,N}$}
	
	\begin{definicao}
		Seja $N\ge 3$ um inteiro. A \textbf{grade} $G_{3,N}$ é uma matriz com $3$ linhas e
		$N$ colunas, cujas células recebem os números ímpares positivos em ordem linha a
		linha. A célula na linha $r\in\{1,2,3\}$ e coluna $c\in\{1,\dots,N\}$ é:
		\[
		f(r,c) = 2\bigl[(r-1)N + c\bigr] - 1.
		\]
	\end{definicao}
	
	As três linhas contêm respectivamente:
	\begin{align*}
		\text{Linha 1:}&\quad 1,\;3,\;\dots,\;2N-1,\\
		\text{Linha 2:}&\quad 2N+1,\;2N+3,\;\dots,\;4N-1,\\
		\text{Linha 3:}&\quad 4N+1,\;4N+3,\;\dots,\;6N-1.
	\end{align*}
	
	\subsection{Janelas $3\times 3$}
	
	Para $i=1,\dots,N-2$, a \textbf{janela} $W_i$ é a submatriz de $G_{3,N}$ formada
	pelas colunas $i$, $i+1$ e $i+2$. Dentro de $W_i$ destacamos três
	\textbf{configurações canônicas}:
	\begin{itemize}
		\item $\mathbf{D_1}$ (diagonal principal): células $(1,i)$, $(2,i+1)$, $(3,i+2)$;
		\item $\mathbf{D_2}$ (diagonal secundária): células $(3,i)$, $(2,i+1)$, $(1,i+2)$;
		\item $\mathbf{C_v}$ (coluna central): células $(1,i+1)$, $(2,i+1)$, $(3,i+1)$.
	\end{itemize}
	
	\begin{figure}[h]
		\centering
		\fbox{\begin{minipage}{0.55\textwidth}
				\centering
				\small
				\begin{tabular}{ccc}
					$f(1,i)$ & $f(1,i+1)$ & $f(1,i+2)$\\[2pt]
					$f(2,i)$ & $f(2,i+1)$ & $f(2,i+2)$\\[2pt]
					$f(3,i)$ & $f(3,i+1)$ & $f(3,i+2)$
				\end{tabular}
				\\[3mm]
				\emph{$D_1$: $(1,i),(2,i+1),(3,i+2)$;\quad
					$D_2$: $(3,i),(2,i+1),(1,i+2)$;\quad
					$C_v$: $(1,i+1),(2,i+1),(3,i+1)$.}
		\end{minipage}}
		\caption{Janela $W_i$ e as três configurações canônicas.}
		\label{fig:janela}
	\end{figure}
	
	\begin{observacao}
		O elemento $f(2,i+1)$ pertence simultaneamente a $D_1$, $D_2$ e $C_v$.
		Ele é o \textit{elemento de interseção} das três configurações e,
		como veremos, o pivô aritmético da janela.
	\end{observacao}
	
	%----------------------------------------------------------------------
	\section{Lema da soma constante}
	\label{sec:lema}
	%----------------------------------------------------------------------
	
	\begin{lema}[Soma Constante]
		\label{lema:soma}
		Para qualquer janela $W_i$ ($1\le i\le N-2$) da grade $G_{3,N}$, as três
		configurações $D_1$, $D_2$ e $C_v$ têm a mesma soma, dada por
		\[
		\boxed{\;D_1 \;=\; D_2 \;=\; C_v \;=\; S(i) \;:=\; 3(2N+2i+1)\;}.
		\]
	\end{lema}
	
	\begin{proof}
		Usando a fórmula $f(r,c)=2[(r-1)N+c]-1$:
		
		\medskip\noindent\textbf{Diagonal $D_1$:}
		\begin{align*}
			f(1,i)   &= 2i-1,\\
			f(2,i+1) &= 2\bigl[N+(i+1)\bigr]-1 = 2N+2i+1,\\
			f(3,i+2) &= 2\bigl[2N+(i+2)\bigr]-1 = 4N+2i+3.
		\end{align*}
		\[
		D_1 = (2i-1)+(2N+2i+1)+(4N+2i+3) = 6N+6i+3 = 3(2N+2i+1).
		\]
		
		\medskip\noindent\textbf{Diagonal $D_2$:}
		\begin{align*}
			f(3,i)   &= 2[2N+i]-1 = 4N+2i-1,\\
			f(2,i+1) &= 2N+2i+1,\\
			f(1,i+2) &= 2(i+2)-1 = 2i+3.
		\end{align*}
		\[
		D_2 = (4N+2i-1)+(2N+2i+1)+(2i+3) = 6N+6i+3.
		\]
		
		\medskip\noindent\textbf{Coluna central $C_v$:}
		\begin{align*}
			f(1,i+1) &= 2(i+1)-1 = 2i+1,\\
			f(2,i+1) &= 2N+2i+1,\\
			f(3,i+1) &= 2[2N+(i+1)]-1 = 4N+2i+1.
		\end{align*}
		\[
		C_v = (2i+1)+(2N+2i+1)+(4N+2i+1) = 6N+6i+3.
		\]
		
		Portanto $D_1=D_2=C_v=6N+6i+3=3(2N+2i+1)$.
	\end{proof}
	
	\begin{observacao}
		A demonstração é puramente algébrica: a igualdade das somas é um fato sobre a
		bijeção $f(r,c)$ e independe de qualquer propriedade de primalidade dos elementos.
	\end{observacao}
	
	%----------------------------------------------------------------------
	\section{Corolário âncora}
	\label{sec:ancora}
	%----------------------------------------------------------------------
	
	O Lema~\ref{lema:soma} tem uma consequência imediata que será central para o programa
	do crivo geométrico.
	
	\begin{corolario}[Invariante Âncora]
		\label{cor:ancora}
		O elemento da linha mediana na coluna central de $W_i$ satisfaz:
		\[
		f(2,\,i+1) \;=\; \frac{S(i)}{3} \;=\; 2N+2i+1.
		\]
		Em particular, $f(2,i+1)$ divide $S(i)$ exatamente e é um \emph{invariante} da
		janela $W_i$: é determinado apenas por $N$ e pela posição $i$, sem qualquer
		informação de primalidade.
	\end{corolario}
	
	\begin{proof}
		Direto de $f(2,i+1)=2[(2-1)N+(i+1)]-1=2N+2i+1$ e $S(i)/3=2N+2i+1$.
	\end{proof}
	
	O Corolário~\ref{cor:ancora} permite definir o \textbf{complemento aritmético} de
	qualquer célula $n$ em relação à janela $W_i$ que a contém.
	
	\begin{definicao}[Complemento oracle-free]
		\label{def:complemento}
		Para $n=f(r,c)$ e janela $W_i$ contendo $n$, define-se:
		\[
		C(n,i) \;:=\; S(i) - n \;=\; 3(2N+2i+1) - n.
		\]
	\end{definicao}
	
	\begin{proposicao}[Paridade do complemento]
		\label{prop:paridade}
		Para todo elemento $n$ da grade $G_{3,N}$ e toda janela $W_i$ contendo $n$,
		o complemento $C(n,i)$ é sempre par.
	\end{proposicao}
	
	\begin{proof}
		Todo elemento $f(r,c)=2[(r-1)N+c]-1$ é ímpar por construção. A soma
		$S(i)=3(2N+2i+1)$ é o produto de $3$ por um número ímpar, logo ímpar.
		Portanto $C(n,i)=S(i)-n=\text{ímpar}-\text{ímpar}$ é par.
	\end{proof}
	
	\begin{observacao}
		A paridade garantida de $C(n,i)$ é essencial para a aplicação em~\cite{bandeira2026c}:
		a Série Singular de Hardy–Littlewood $\mathfrak{S}(C)=\prod_{p\mid C,\, p>2}
		\frac{p-1}{p-2}$ é definida sobre inteiros pares, e a
		Proposição~\ref{prop:paridade} garante que $C(n,i)$ satisfaz essa hipótese
		para toda célula da grade, tornando $\mathfrak{S}(C(n,i))$ bem-definida sem
		qualquer restrição adicional.
	\end{observacao}
	
	\begin{observacao}
		A denominação ``oracle-free'' refere-se ao fato de que $C(n,i)=3(2N+2i+1)-n$
		depende apenas de $N$, da posição $i$ e do valor $n$ — nenhum dos quais requer
		um teste de primalidade. Em contraste, o peso $\Omega_N(n)$ de~\cite{bandeira2026}
		requer saber se cada vizinho de $n$ é primo. O complemento $C(n,i)$ é o objeto
		que permite contornar essa dependência.
	\end{observacao}
	
	%----------------------------------------------------------------------
	\section{Estrutura aritmética das configurações}
	\label{sec:aritm}
	%----------------------------------------------------------------------
	
	A coincidência das três somas não é acidental: cada configuração canônica possui uma
	estrutura aritmética interna precisa, que revelamos nesta seção.
	
	\subsection{As configurações são progressões aritméticas}
	
	\begin{proposicao}[Progressões aritméticas]
		\label{prop:pa}
		Cada uma das três configurações canônicas forma uma progressão aritmética:
		\begin{align*}
			D_1 &:\quad 2i-1,\quad 2N+2i+1,\quad 4N+2i+3
			\quad(\text{razão }\ r_{D_1}=2N+2),\\
			D_2 &:\quad 2i+3,\quad 2N+2i+1,\quad 4N+2i-1
			\quad(\text{razão }\ r_{D_2}=2N-2),\\
			C_v &:\quad 2i+1,\quad 2N+2i+1,\quad 4N+2i+1
			\quad(\text{razão }\ r_{C_v}=2N).
		\end{align*}
	\end{proposicao}
	
	\begin{proof}
		As expressões dos elementos foram calculadas na demonstração do Lema~\ref{lema:soma}.
		As diferenças entre termos consecutivos são:
		\begin{align*}
			D_1:\quad & (2N+2i+1)-(2i-1) = 2N+2,\quad (4N+2i+3)-(2N+2i+1)=2N+2.\\
			D_2:\quad & \text{Em ordem crescente: } f(1,i+2)< f(2,i+1)< f(3,i)\\
			& (2N+2i+1)-(2i+3) = 2N-2,\quad (4N+2i-1)-(2N+2i+1)=2N-2.\\
			C_v:\quad & (2N+2i+1)-(2i+1) = 2N,\quad (4N+2i+1)-(2N+2i+1)=2N.
		\end{align*}
		Em cada caso a diferença entre termos consecutivos é constante.
	\end{proof}
	
	\begin{observacao}
		As três razões $2N+2$, $2N-2$ e $2N$ são simétricas em torno de $2N$ e somam
		$6N$. As razões dependem exclusivamente de $N$: janelas distintas (valores diferentes
		de $i$) produzem progressões com as mesmas razões mas termos iniciais distintos.
	\end{observacao}
	
	\subsection{Diferenças constantes entre linhas}
	
	\begin{proposicao}[Diferenças entre linhas]
		\label{prop:dif}
		Para qualquer coluna $c$, a diferença entre elementos de linhas consecutivas é
		constante:
		\[
		f(r+1,c) - f(r,c) = 2N \quad \text{para todo } r\in\{1,2\},\ c\in\{1,\dots,N\}.
		\]
		Consequentemente, a grade $G_{3,N}$ é \emph{verticalmente uniforme}: cada coluna
		é uma progressão aritmética de razão $2N$.
	\end{proposicao}
	
	\begin{proof}
		$f(r+1,c)-f(r,c)=2[(r)N+c]-1-\bigl(2[(r-1)N+c]-1\bigr)=2N$.
	\end{proof}
	
	Em particular, dentro de qualquer janela $W_i$, as linhas 1, 2 e 3 na mesma coluna
	$j\in\{i,i+1,i+2\}$ satisfazem $f(2,j)-f(1,j)=f(3,j)-f(2,j)=2N$, ou seja,
	cada coluna da janela é também uma progressão aritmética de razão $2N$ — a mesma
	razão de $C_v$.
	
	\subsection{Estrutura modular dos elementos}
	
	\begin{proposicao}[Uniformidade modular]
		\label{prop:mod}
		Para todo $r\in\{1,2,3\}$ e $c\in\{1,\dots,N\}$:
		\[
		f(r,c) \;\equiv\; 2c - 1 \pmod{2N}.
		\]
		Em particular, todos os elementos de uma mesma coluna têm o mesmo resíduo módulo
		$2N$, independentemente da linha.
	\end{proposicao}
	
	\begin{proof}
		$f(r,c)=2(r-1)N+2c-1\equiv 2c-1\pmod{2N}$.
	\end{proof}
	
	\begin{corolario}
		Os elementos das três configurações canônicas têm os seguintes resíduos módulo $2N$:
		\begin{align*}
			D_1 &:\quad (2i-1),\;(2i+1),\;(2i+3)\pmod{2N},\\
			D_2 &:\quad (2i-1),\;(2i+1),\;(2i+3)\pmod{2N},\\
			C_v &:\quad (2i+1),\;(2i+1),\;(2i+1)\pmod{2N}.
		\end{align*}
		As diagonais $D_1$ e $D_2$ têm os mesmos três resíduos (em ordens diferentes), e
		os três elementos de $C_v$ são congruentes entre si módulo $2N$.
	\end{corolario}
	
	\begin{proof}
		Aplica-se diretamente a Proposição~\ref{prop:mod} às colunas $i$, $i+1$, $i+2$.
		Para $C_v$, os três elementos estão na mesma coluna $i+1$, logo todos têm
		resíduo $2(i+1)-1=2i+1$.
	\end{proof}
	
	\begin{observacao}
		A uniformidade modular de $C_v$ implica que se $2N$ tem fatores primos pequenos
		(i.e., $N$ é altamente composto), então todos os três elementos de $C_v$ ficam na
		mesma classe de resíduo modulo esses fatores — criando obstruções à primalidade
		simultânea. Esse fenômeno é a origem das ``obstruções modulares'' descritas
		em~\cite{bandeira2026} para $N$ com fatoração em $\{2,3,5\}$, e é precisamente
		o problema que $\Omega^*_N$ contorna ao trabalhar com o complemento $C(n,i)$ em
		vez dos próprios elementos.
	\end{observacao}
	
	%----------------------------------------------------------------------
	\section{Por que $G_{3,N}$? O caso privilegiado entre as grades $G_{k,N}$}
	\label{sec:general}
	%----------------------------------------------------------------------
	
	A escolha de $G_{3,N}$ como objeto central desta série não é arbitrária nem
	meramente convencional. Esta seção responde à pergunta natural: \emph{por que estudar $G_{3,N}$
		e não $G_{k,N}$ para $k$ arbitrário?} A resposta é que $k=3$ é o único
	valor que satisfaz simultaneamente três condições necessárias para o programa
	do crivo geométrico — e essa coincidência é o que torna os resultados das
	seções anteriores possíveis.
	
	As três condições são:
	\begin{enumerate}
		\item \textbf{Existência da âncora inteira.} Para que $S(i)/k$ seja um elemento
		real da grade (e não apenas um valor racional), é preciso que $k$ seja
		ímpar — caso contrário não há linha mediana inteira.
		\item \textbf{Complemento $C(n,i)$ sempre par.} A Série Singular
		$\mathfrak{S}(C)$ de Hardy–Littlewood requer que $C$ seja par;
		isso ocorre quando $S_k(i)$ é ímpar, o que exige $k$ ímpar e o
		elemento âncora ímpar — ambos garantidos para $k=3$.
		\item \textbf{Igualdade tripla $D_1=D_2=C_v=S(i)$.} Essa propriedade é
		específica da geometria de janelas $3\times 3$: com $k=3$ linhas e
		$3$ colunas, as duas diagonais e a coluna central são os únicos
		percursos que cobrem uma célula por linha e passam pelo pivô
		$f(2,i+1)$. Para $k>3$, janelas mais altas admitem percursos
		diagonais adicionais que \emph{não} têm a mesma soma, destruindo
		a unicidade do Lema~\ref{lema:soma}.
	\end{enumerate}
	
	O caso $k=3$ é o \emph{menor} valor ímpar para o qual as três condições
	coexistem de forma não-trivial. Para $k=1$ a grade degenera numa linha e
	não há configurações diagonais. Para $k=3$, as condições se alinham
	precisamente — e é isso que faz de $S(i)=3(2N+2i+1)$, do Invariante
	Âncora $f(2,i+1)=S(i)/3$ e do complemento $C(n,i)=S(i)-n$ (sempre par)
	objetos bem definidos e aritmeticamente úteis.
	
	A análise a seguir torna esse argumento preciso, mostrando o que se perde
	para cada $k\ne 3$.
	
	\begin{definicao}[Grade $G_{k,N}$]
		Para $k,N\ge 2$ inteiros, a grade $G_{k,N}$ é a matriz $k\times N$ cujas células
		recebem os primeiros $kN$ números ímpares linha a linha:
		\[
		f_k(r,c) = 2\bigl[(r-1)N+c\bigr]-1,\quad r\in\{1,\dots,k\},\;c\in\{1,\dots,N\}.
		\]
	\end{definicao}
	
	\subsection{Soma da coluna central}
	
	Para qualquer coluna $j\in\{1,\dots,N\}$, a soma dos $k$ elementos da coluna é:
	\begin{align*}
		\Sigma_k(j) &:= \sum_{r=1}^{k} f_k(r,j)
		= \sum_{r=1}^{k} \bigl[2(r-1)N + 2j -1\bigr]\\
		&= k(2j-1) + 2N\sum_{r=1}^{k}(r-1)
		= k(2j-1) + 2N\cdot\frac{k(k-1)}{2}\\
		&= k\bigl[(k-1)N + 2j - 1\bigr].
	\end{align*}
	
	\begin{proposicao}[$k$ ímpar: existência da âncora]
		\label{prop:kpar}
		Quando $k$ é ímpar, existe uma linha mediana $r^*=\frac{k+1}{2}$ (inteira), e
		\[
		\Sigma_k(j) = k\cdot f_k(r^*,j).
		\]
		O elemento $f_k(r^*,j)$ é o \textbf{invariante âncora} da coluna $j$ em
		$G_{k,N}$.
	\end{proposicao}
	
	\begin{proof}
		$f_k(r^*,j)=2\bigl[(r^*-1)N+j\bigr]-1=2\bigl[\frac{k-1}{2}N+j\bigr]-1=(k-1)N+2j-1$.
		Portanto $k\cdot f_k(r^*,j)=k[(k-1)N+2j-1]=\Sigma_k(j)$.
	\end{proof}
	
	\begin{observacao}[$k=3$]
		Para $k=3$, $r^*=2$ e $\Sigma_3(j)=3(2N+2j-1)=3f_3(2,j)=S(j-1)$ (com $j=i+1$),
		recuperando o Corolário~\ref{cor:ancora}.
	\end{observacao}
	
	\begin{proposicao}[$k=2$: ausência de âncora]
		\label{prop:k2}
		Para $k=2$, a soma da coluna é
		\[
		\Sigma_2(j) = 2(N+2j-1),
		\]
		que é sempre par. Não existe linha mediana inteira, logo não há âncora no sentido
		do Corolário~\ref{cor:ancora}. O complemento $C_2(n,j):=\Sigma_2(j)-n$ é ímpar
		(par menos ímpar), de modo que a Série Singular de Hardy–Littlewood não se aplica
		diretamente.
	\end{proposicao}
	
	\begin{proof}
		$\Sigma_2(j)=2[(2j-1)+N]= 2(N+2j-1)$. Para $n=f_2(r,j)$ ímpar,
		$C_2(n,j)=2(N+2j-1)-n$ é par menos ímpar, portanto ímpar.
	\end{proof}
	
	\begin{proposicao}[$k=4$: âncora racional e complemento par]
		\label{prop:k4}
		Para $k=4$, a soma da coluna é
		\[
		\Sigma_4(j) = 4(3N+2j-1),
		\]
		sempre divisível por $4$. Não há linha mediana inteira (a ``âncora'' seria
		$3N+2j-1$, que não corresponde a $f_4(r,j)$ para nenhum $r$ inteiro em geral).
		Entretanto, o complemento $C_4(n,j):=\Sigma_4(j)-n$ é \emph{sempre par}
		(par menos ímpar é ímpar; mas $\Sigma_4(j)=4(\cdots)$ é par, logo $C_4=\text{par}-\text{ímpar}=\text{ímpar}$).
	\end{proposicao}
	
	\begin{proof}
		$\Sigma_4(j)=4[(4-1)N+2j-1]=4(3N+2j-1)$. Este é divisível por $4$, logo par.
		Para $n$ ímpar, $C_4=4(3N+2j-1)-n=\text{par}-\text{ímpar}=\text{ímpar}$.
	\end{proof}
	
	\begin{observacao}
		O caso $k=4$ difere do caso $k=3$ em dois aspectos: (i) não há âncora inteira
		(a linha mediana cairia entre as linhas 2 e 3); (ii) o complemento $C_4$ é ímpar,
		não par, de modo que $\mathfrak{S}(C_4)$ não é diretamente aplicável. Uma variante
		da Série Singular para inteiros ímpares pode ser definida, mas requer adaptações.
		Para $k=5$, a âncora existe ($r^*=3$) e o complemento é par, recuperando as
		propriedades favoráveis de $k=3$.
	\end{observacao}
	
	\subsection{Tabela resumo}
	
	\begin{center}
		\begin{tabular}{ccccl}
			\hline
			$k$ & $\Sigma_k(j)$ & Âncora inteira? & $C_k(n,j)$ par? & Observação\\
			\hline
			$2$ & $2(N+2j-1)$ & Não & Não (ímpar) & Sem âncora nem $\mathfrak{S}$\\
			$3$ & $3(2N+2j-1)$ & Sim ($r^*=2$) & Sim & Caso principal\\
			$4$ & $4(3N+2j-1)$ & Não & Não (ímpar) & $\mathfrak{S}$ não se aplica\\
			$5$ & $5(4N+2j-1)$ & Sim ($r^*=3$) & Sim & Análogo a $k=3$\\
			$k$ ímpar & $k[(k-1)N+2j-1]$ & Sim ($r^*=\frac{k+1}{2}$) & Sim & Caso geral favorável\\
			$k$ par & $k[\frac{(k-1)}{1}N+\cdots]$ & Não & Não & Requer tratamento distinto\\
			\hline
		\end{tabular}
	\end{center}
	
	\begin{observacao}[Unicidade de $k=3$ para o programa do crivo geométrico]
		A tabela mostra que $k=3$ e $k=5$ são os primeiros casos com âncora inteira
		e complemento par. Entretanto, $k=5$ falha na condição~3: janelas $5\times 3$
		admitem dez percursos diagonais distintos entre linhas 1 e 5, e apenas dois
		deles (as diagonais extremas) passam pelo elemento central; a coluna central
		tem cinco elementos. A igualdade tripla $D_1=D_2=C_v$ --- que é o
		Lema~\ref{lema:soma} e o fundamento de tudo que se segue --- não tem análogo
		direto para $k=5$.
		
		Portanto, \textbf{$G_{3,N}$ é o caso natural mínimo} onde as três condições
		coexistem: âncora inteira ($r^*=2$), complemento par (logo $\mathfrak{S}(C)$
		bem-definida), e igualdade das três configurações canônicas (logo $S(i)$ único
		e bem-definido). É por isso que toda a série --- o peso $\Omega_N$
		de~\cite{bandeira2026}, o Invariante Âncora desta seção e o estimador
		$\Omega^*_N$ de~\cite{bandeira2026c} --- é desenvolvida exclusivamente em
		$G_{3,N}$.
	\end{observacao}
	
	%----------------------------------------------------------------------
	\section{Conexão com as conjecturas de Goldbach}
	\label{sec:goldbach}
	%----------------------------------------------------------------------
	
	\subsection{Configurações totalmente primas e a Conjectura forte}
	
	Dizemos que uma configuração canônica é \textbf{totalmente prima} quando todos os
	seus três elementos são primos.
	
	Se a configuração $D_1$ é totalmente prima, com
	$f(1,i)=p_a$, $f(2,i+1)=p_b$, $f(3,i+2)=p_c$ primos, então os três pares produzem
	números pares:
	\begin{align*}
		p_a + p_b &= (2i-1)+(2N+2i+1) = 2(N+2i),\\
		p_a + p_c &= (2i-1)+(4N+2i+3) = 2(2N+2i+1),\\
		p_b + p_c &= (2N+2i+1)+(4N+2i+3) = 2(3N+2i+2).
	\end{align*}
	Cada uma é um número par $>2$ que é soma de dois primos: três instâncias da
	Conjectura de Goldbach forte. O mesmo raciocínio aplica-se a $D_2$ e $C_v$,
	cada uma gerando três representações com pares distintos.
	
	\begin{observacao}
		Em particular, o par $p_a+p_c=2(2N+2i+1)=2S(i)/3$ é sempre o dobro da âncora
		$f(2,i+1)$. Assim, o elemento central da janela não é apenas um invariante da soma:
		ele é o semi-valor de um dos pares de Goldbach testemunhados pela configuração.
	\end{observacao}
	
	\subsection{Conjectura fraca (ternária)}
	
	Pelo Lema~\ref{lema:soma}, se $D_1$ é totalmente prima:
	\[
	p_a + p_b + p_c = S(i) = 3(2N+2i+1),
	\]
	que é um ímpar $>5$ para $N\ge 3$. Logo, cada trio primo na grade fornece uma
	\textbf{representação da Conjectura de Goldbach ternária} para o ímpar $S(i)$.
	
	\begin{observacao}
		Embora a Conjectura fraca seja um teorema~\cite{helfgott2013}, a estrutura
		geométrica aqui apresentada exibe uma família explícita de representações ternárias
		parametrizada por $N$ e $i$, com a propriedade de que cada trio primo também
		testemunha três representações binárias. Isso unifica as duas formas da conjectura
		num único objeto geométrico.
	\end{observacao}
	
	\subsection{O complemento $C(n,i)$ como fundamento oracle-free}
	
	A Definição~\ref{def:complemento} e a Proposição~\ref{prop:paridade} fornecem o
	elemento central para o programa do crivo geométrico. Dado um par $a=2M$ que se
	deseja representar como soma de dois primos, e dado um primo $p<a/2$ que pertence
	a $G_{3,N}$ como célula $n=p$ na janela $W_i$, o complemento
	\[
	C(p,i) = S(i) - p
	\]
	é um número par determinístico, computável a partir de $N$, $i$ e $p$, sem
	nenhum oráculo de primalidade. Se $C(p,i)$ for primo, então $p+C(p,i)=S(i)$
	é uma representação ternária; mas mais relevante para Goldbach binária, a
	estrutura aditiva de $C(p,i)$ — medida pela Série Singular $\mathfrak{S}(C(p,i))$
	— estima quão ``Goldbach-favorável'' é o ambiente aritmético de $p$ em $W_i$.
	
	Em~\cite{bandeira2026c}, esse princípio é formalizado no estimador
	\[
	\Omega^*_N(n) := \frac{1}{|W(n)|}\sum_{i\in W(n)} \mathfrak{S}\bigl(C(n,i)\bigr),
	\]
	onde $W(n)$ é o conjunto de janelas contendo $n$ e
	$\mathfrak{S}(C)=\prod_{p\mid C,\,p>2}\frac{p-1}{p-2}$.
	A verificação computacional em~\cite{bandeira2026c} mostra que $\Omega^*_N$
	discrimina primos de semiprimos em todo $N$ testado, incluindo os $N$ com obstruções
	modulares onde $\Omega_N$ falha.
	
	\begin{observacao}[Cada $D_1$ primo é uma célula de $\bigcup A_p$]
		Denotando por $A_p$ o conjunto dos pares $(p,q)$ com $p+q=2M$ e $p,q$ primos,
		cada configuração totalmente prima $D_1=(p_a,p_b,p_c)$ contribui com três elementos
		$\{(p_a,p_b),(p_a,p_c),(p_b,p_c)\}$ para $\bigcup_{2M} A_{2M}$.
		O Lema~\ref{lema:soma} garante que os três pares somam valores \emph{distintos}
		($2(N+2i)$, $2(2N+2i+1)$, $2(3N+2i+2)$), de modo que uma única configuração prima
		povoa três instâncias distintas de $A_p$.
	\end{observacao}
	
	%----------------------------------------------------------------------
	\section{Exemplos numéricos}
	\label{sec:exemplos}
	%----------------------------------------------------------------------
	
	\subsection{Verificação básica: $N=5$, $i=2$}
	
	A grade $G_{3,5}$ é:
	\[
	\begin{pmatrix}
		1  &  3 &  5 &  7 &  9\\
		11 & 13 & 15 & 17 & 19\\
		21 & 23 & 25 & 27 & 29
	\end{pmatrix}.
	\]
	A janela $W_2$ (colunas 2, 3, 4) é:
	\[
	\begin{pmatrix}
		3  &  5 &  7\\
		13 & 15 & 17\\
		23 & 25 & 27
	\end{pmatrix}.
	\]
	Configurações: $D_1=3+15+27=45$, $D_2=23+15+7=45$, $C_v=5+15+25=45$.
	De fato, $S(2)=3(10+4+1)=45$. A âncora é $f(2,3)=15=45/3$. Nenhuma
	configuração é totalmente prima (15 e 25 não são primos), mas a identidade algébrica
	permanece. O complemento de $n=3$ na janela $W_2$: $C(3,2)=45-3=42$ (par).
	
	\subsection{Trio primo e representações de Goldbach: $N=20$, $i=3$}
	
	\begin{align*}
		f(1,3)   &= 2\cdot 3 - 1 = 5,\\
		f(2,4)   &= 2[20+4]-1 = 47,\\
		f(3,5)   &= 2[40+5]-1 = 89.
	\end{align*}
	Os três são primos. Âncora: $f(2,4)=47=S(3)/3=141/3$. Verificação:
	$D_1=5+47+89=141=3\cdot 47$.\checkmark
	
	\medskip\noindent\textbf{Representações binárias (Goldbach forte):}
	\[
	52 = 5+47,\qquad 94 = 5+89,\qquad 136 = 47+89.
	\]
	\textbf{Representação ternária (Goldbach fraca):}
	\[
	141 = 5+47+89,\quad 141\text{ ímpar}>5.\;\checkmark
	\]
	
	Complementos oracle-free de cada elemento em $W_3$:
	$C(5,3)=141-5=136$ (par), $C(47,3)=94$ (par), $C(89,3)=52$ (par).
	
	\subsection{Trio primo em $D_2$: $N=10$, $i=4$}
	
	\begin{align*}
		f(3,4)   &= 2[20+4]-1 = 47,\\
		f(2,5)   &= 2[10+5]-1 = 29,\\
		f(1,6)   &= 2\cdot 6 - 1 = 11.
	\end{align*}
	Todos primos. $S(4)=3(20+8+1)=87$. Âncora: $f(2,5)=29=87/3$.\checkmark
	
	Representações binárias:
	\[
	76=47+29,\qquad 58=47+11,\qquad 40=29+11.
	\]
	Representação ternária: $87=47+29+11$.\checkmark
	
	\subsection{Exemplo com a coluna central $C_v$: $N=15$, $i=5$}
	
	\begin{align*}
		f(1,6) &= 11,\quad f(2,6)=2[15+6]-1=41,\quad f(3,6)=2[30+6]-1=71.
	\end{align*}
	Todos primos. $S(5)=3(30+10+1)=123$. Âncora: $41=123/3$.\checkmark
	
	$C_v=11+41+71=123$. Representações binárias: $52=11+41$, $82=11+71$,
	$112=41+71$. Representação ternária: $123=11+41+71$.\checkmark
	
	\subsection{Obstrução modular e o papel de $C(n,i)$}
	
	Para $N=1800=2^3\cdot 3^2\cdot 5^2$, a grade tem fortes restrições modulares:
	cada coluna $c$ tem $f(r,c)\equiv 2c-1\pmod{2N}$, e como $2N=3600$ é divisível
	por $2$, $3$, $4$, $5$, $6$, $9$, $10$, $12$, $15$, $16$, $18$, $20$, $24$, $25$,
	$36$, $\dots$, muitos resíduos ficam obstruídos para primalidade simultânea.
	Isso faz $\Omega_N$ perder poder discriminatório (ver~\cite{bandeira2026c},
	Tabela~1, $N=1800$: $p$-valor KS $=0{,}118$, não significativo).
	
	Entretanto, $C(n,i)=S(i)-n$ tem fatoração determinada pela \emph{posição} de $n$,
	não pelas obstruções modulares de $N$. Para $N=1800$, o estimador $\Omega^*_N$
	obtém $p$-valor KS $=3{,}15\times 10^{-47}$, discriminando primos de semiprimos
	com altíssima significância estatística — confirmando que o complemento $C$ captura
	um sinal aritmético independente das obstruções.
	
	%----------------------------------------------------------------------
	\section{Conclusão}
	\label{sec:conclusao}
	%----------------------------------------------------------------------
	
	Apresentamos uma propriedade aritmética elementar da grade $3\times N$ — a
	igualdade das somas das configurações canônicas $D_1$, $D_2$ e $C_v$ em toda
	janela $W_i$ — e desenvolvemos suas consequências em quatro direções.
	
	\medskip
	
	A primeira é o \textbf{Invariante Âncora}: $f(2,i+1)=S(i)/3$ é um elemento
	determinístico da grade que particiona $S(i)$ em terços, e cujo conhecimento
	basta para calcular o complemento $C(n,i)=S(i)-n$ de qualquer célula sem oráculo
	de primalidade. Essa é a pedra angular do trabalho~\cite{bandeira2026c}, onde
	$C(n,i)$ fundamenta o estimador $\Omega^*_N$ que resolve o obstáculo computacional
	central do programa do crivo geométrico.
	
	\medskip
	
	A segunda é a \textbf{estrutura aritmética interna}: cada configuração canônica é
	uma progressão aritmética com razão que depende linearmente de $N$ ($2N+2$ para
	$D_1$, $2N-2$ para $D_2$, $2N$ para $C_v$), e todos os elementos de uma mesma
	coluna são congruentes módulo $2N$. Essa uniformidade modular explica as obstruções
	que enfraquecem $\Omega_N$ para $N$ altamente composto, e mostra por que o
	complemento $C$ escapa dessas obstruções.
	
	\medskip
	
	A terceira é a \textbf{generalização para $k\times N$}: para $k$ ímpar existe sempre
	uma linha mediana inteira $r^*=(k+1)/2$ que é âncora da janela, e o complemento
	$C_k$ é sempre par — propriedades que tornam $\mathfrak{S}(C_k)$ bem-definida.
	Para $k$ par não há âncora inteira e o complemento é ímpar, o que exige adaptações.
	
	\medskip
	
	A quarta é a \textbf{conexão com Goldbach}: cada trio primo na grade produz
	simultaneamente três representações binárias (com pares distintos) e uma
	representação ternária (para o ímpar $S(i)$). Assim, a identidade algébrica do
	Lema~\ref{lema:soma} transforma cada trio primo em testemunha múltipla de ambas
	as formas da conjectura.
	
	\medskip
	
	Os três artigos da série formam um framework coerente: \cite{bandeira2026}
	introduz $G_{3,N}$, o peso $\Omega_N$ e o argumento probabilístico
	de Borel–Cantelli; o presente artigo fornece a base algébrica — a soma constante,
	a âncora e o complemento $C$; e \cite{bandeira2026c} usa essa base para construir
	$\Omega^*_N$ e estabelecer a discriminação oracle-free entre primos e semiprimos.
	
	\begin{thebibliography}{9}
		\bibitem{bandeira2026}
		T. Bandeira,
		``Uma Abordagem Unificada para a Conjectura de Goldbach: Estrutura Combinatória,
		Saturação Geométrica e o Elo entre Trios e Pares Primos'',
		Preprint, maio de 2026.
		
		\bibitem{bandeira2026c}
		T. Bandeira,
		``O Invariante Âncora da Grade $G_{3,N}$ e um Estimador Geométrico Oracle-Free
		para a Conjectura de Goldbach'',
		Preprint, maio de 2026.
		
		\bibitem{helfgott2013}
		H. A. Helfgott,
		``Major arcs for Goldbach's theorem'',
		\textit{arXiv}:1305.2897, 2013.
		
		\bibitem{hardylittlewood1923}
		G. H. Hardy e J. E. Littlewood,
		``Some problems of `Partitio Numerorum'; III: On the expression of a number
		as a sum of primes'',
		\textit{Acta Mathematica}, vol.~44, pp.~1--70, 1923.
		
		\bibitem{selberg1947}
		A. Selberg,
		``On an elementary method in the theory of primes'',
		\textit{Norske Vid. Selsk. Forh.} (Trondheim), vol.~19, no.~18, pp.~64--67, 1947.
	\end{thebibliography}
	
\end{document}